\documentclass[12pt]{article}

% ---------- Basic setup ----------
\usepackage[margin=1in]{geometry}
\usepackage[T1]{fontenc}
\usepackage{lmodern}
\usepackage{microtype}
\usepackage{setspace}
\usepackage{algorithm}
\usepackage{algpseudocode}
\usepackage[most]{tcolorbox}
\usepackage{algorithmic}


\onehalfspacing

% Paragraph look: clear, readable
\setlength{\parindent}{12pt}
\setlength{\parskip}{0.4em}

% ---------- Math & symbols ----------
\usepackage{amsmath, amssymb, amsthm}
\usepackage{siunitx} % \SI{}{} and nice numbers
\sisetup{round-mode=places, round-precision=2, detect-all}

% ---------- Tables & figures ----------
\usepackage{graphicx}
\graphicspath{{./}} % where your figures live
\usepackage{booktabs}
\usepackage[labelfont=bf,font=small]{caption}
\usepackage{subcaption}
\usepackage{float}      % [H] exact float placement
\usepackage{placeins}   % \FloatBarrier to stop floats crossing sections

% ---------- Links & clever refs ----------
\usepackage[hidelinks]{hyperref}
\usepackage[capitalise,nameinlink]{cleveref}

% ---------- Author blocks ----------
\usepackage{authblk}

% ---------- Bibliography ----------
\usepackage[round,authoryear]{natbib}
\bibliographystyle{apalike}

% ---------- Custom commands ----------
\newcommand{\VaR}{\mathrm{VaR}}
\newcommand{\ES}{\mathrm{ES}}
\newcommand{\Prob}{\mathbb{P}}
\newcommand{\E}{\mathbb{E}}
\newcommand{\R}{\mathbb{R}}
\newcommand{\degC}{\,^\circ\mathrm{C}}

% ---------- Title ----------
\title{Quantifying Heat-Hazard Risk with Spatial Clustering }
\author[1,2]{Widad Kchikach}
\author[2]{Thorsten Dickhaus}
\author[1,3]{Gerrit Lohmann}
\affil[1]{Alfred Wegener Institute Helmholtz Centre for Polar and Marine Research, Bremerhaven, Germany}
\affil[2]{Institute for Statistics, University of Bremen, Bremen, Germany}
\affil[3]{Department of Environmental Physics, University of Bremen, Bremen, Germany}
\date{\today}

\begin{document}
\maketitle

% (Removed the footnote per your request.


\begin{abstract}
We develop a global framework for quantifying heat-hazard risk by combining annual
maximum temperatures (1981--2020), present-day population exposure, and a
data-driven regionalisation of extreme-temperature behaviour. To identify regions where
extreme heat events tend to occur simultaneously, we fit a spatial extremes model
(max-stable process) to the annual maxima and derive clusters based on the estimated
extremal dependence. For each grid cell, we define a locally adaptive temperature
threshold, compute the associated exceedance losses weighted by population exposure,
and summarise tail behaviour using Value-at-Risk (VaR) and Expected Shortfall (ES) at
the 95\% level. Aggregating losses within dependence-consistent clusters allows us to
assess regional exposure to rare but impactful heat events. We present global maps of
cell-wise and cluster-level VaR and ES, discuss key modelling choices (e.g., threshold
selection, static exposure), and outline extensions such as continuous severity
functions, time-varying population, and alternative dependence structures.

\end{abstract}

\noindent\textbf{Keywords:} heat extremes; spatial dependence; max-stable processes; expected shortfall; population exposure; climate risk.



\section{Introduction}

Extreme heat events pose a growing threat to human health, infrastructure, and
socio-economic systems worldwide. Recent decades have witnessed an increase in the
frequency, intensity, and spatial extent of heatwaves, with disproportionate impacts in
regions of high population exposure. Understanding where extreme heat is most likely to
occur, how strongly nearby locations co-experience extreme temperatures, and how many
people are exposed when such events occur is therefore essential for climate-risk
assessment and adaptation planning.

Many hazard-focused studies analyse extreme temperatures locally, typically by fitting
statistical extreme-value models to temperature time series at individual locations.
While this provides insight into marginal behaviour, extreme heat events are inherently
spatial phenomena: high temperatures tend to occur over large regions due to persistent
synoptic-scale atmospheric patterns. Ignoring this spatial structure can lead to
underestimation or mischaracterisation of risk when aggregating losses across regions.

Spatial extremes models address this challenge by jointly representing the behaviour of
extreme temperatures across multiple locations. In particular, max-stable processes
provide a flexible class of models that describe the limiting behaviour of block maxima
and account for the co-occurrence of extremes. These models yield pairwise measures of
extremal dependence, known as extremal coefficients, which quantify how likely two
locations are to experience extreme events in the same years. Such dependence
information provides a natural basis for constructing spatial regions in which extreme
heat tends to occur simultaneously.

The goal of this study is to integrate these ideas into a comprehensive, global
heat-hazard risk pipeline. Our framework combines (i) annual maximum temperatures from
1981--2020, (ii) present-day population exposure, and (iii) spatial clusters derived from
the extremal dependence structure of a fitted max-stable model. This regionalisation
ensures that risk assessments reflect the underlying dependence structure rather than
treating locations as independent.

For each grid cell, we define a locally adaptive temperature threshold based on the
empirical distribution of annual maxima. Losses are computed as exceedances of this
threshold weighted by population exposure, and tail behaviour is summarised using two
standard risk measures: Value-at-Risk (VaR) and Expected Shortfall (ES). These
quantities quantify the severity of rare, high-impact years while retaining clear
interpretability.

Aggregating losses within the dependence-based clusters provides regional risk
assessments that respect the joint behaviour of extreme heat across space. This is
important because the total population exposed during extreme events depends not only on
local temperature distributions but also on how extremes propagate spatially.

The main contributions of this work are:
\begin{enumerate}
    \item A global, dependence-aware framework for quantifying heat-hazard risk that combines spatial extremes modelling, population exposure, and threshold-based loss functions.
    \item A regionalisation of the global climate grid into clusters that share similar patterns of extremal dependence, derived from max-stable modelling of annual temperature maxima.
    \item Cell-wise and cluster-level estimates of Value-at-Risk and Expected Shortfall at the 95\% level, providing interpretable measures of extreme population exposure.
    \item A discussion of key modelling assumptions (e.g., choice of threshold, static exposure) and potential extensions including continuous severity functions, time-varying exposure, and alternative dependence structures.
\end{enumerate}

The remainder of this manuscript is structured as follows. Section (explain each section elements and what it covers)




\section{Mathematical Framework for Spatial Extremes and Heat-Risk Aggregation}

Quantifying heat-related risk requires a statistical framework that captures (i) the behaviour
of extreme temperatures at each grid cell, (ii) the spatial dependence of these extremes,
and (iii) a loss aggregation procedure that respects this dependence. This section presents
the necessary components of extreme-value theory, copulas, extremal coefficients, and the
STANiEXT framework used to construct dependence-aware spatial clusters.

\subsection{Marginal Modelling of Annual Maxima}

Let $i \in \{1,\dots,d\}$ denote a grid cell and $y \in \{1981,\dots,2020\}$ denote the year.
The annual maximum temperature at cell $i$ is defined as
\[
T_{y,i} = \max_{\text{day}} \{ \mathrm{Tmax}_{\text{day},i} \}.
\]
Under classical extreme-value theory, the distribution of block maxima converges to the
Generalized Extreme Value (GEV) family \citep{coles2001introduction}. Hence we model
\[
T_{y,i} \sim \mathrm{GEV}(\mu_i,\sigma_i,\xi_i),
\]
where $\mu_i$ (location), $\sigma_i>0$ (scale), and $\xi_i$ (shape) determine the local
tail behaviour of extremes at each site.

\subsection{Dependence Structures and Extreme-Value Copulas}

In multivariate modelling, copulas allow a joint distribution to be decomposed into its
marginals and a dependence component \citep{nelsen2006introduction}. For variables
$X_1,\dots,X_d$ with continuous marginals $F_1,\dots,F_d$, Sklar’s theorem states:
\[
F(x_1,\dots,x_d) = C\!\left(F_1(x_1),\dots,F_d(x_d)\right),
\]
where $C$ is the copula describing all dependence.

For block maxima, standard copulas are not suitable because they do not capture joint
tail behaviour. Instead, spatial extremes are characterised by \emph{extreme-value copulas},
which arise as limits of componentwise maxima and are stable under the maximum
operation \citep{davison2012statistical}. These copulas capture the co-occurrence of rare,
high-impact temperature events across space.

\subsection{Extremal Coefficients and Tail Dependence}

To analyse tail dependence, annual maxima are standardised to a unit-Fr\'echet
distribution with cumulative distribution function
\[
F_{\mathrm{Fr}}(z) = \exp\!\left(-\frac{1}{z}\right), \qquad z>0.
\]
This monotone transformation preserves dependence.

Let $Z_{y,i}$ and $Z_{y,j}$ denote the standardised maxima at sites $i$ and $j$.
Their joint tail satisfies
\[
\mathbb{P}(Z_{y,i} \le z,\; Z_{y,j} \le z)
=
\exp\!\left(-\frac{\theta_{i,j}}{z}\right), \qquad z>0,
\]
where $\theta_{i,j} \in [1,2]$ is the \emph{extremal coefficient}. It summarises the strength
of extremal dependence:
\[
\theta_{i,j}=1 \iff \text{perfect dependence}, \qquad 
\theta_{i,j}=2 \iff \text{independence}.
\]
The collection of extremal coefficients $\{\theta_{i,j}\}$ across all pairs determines the
underlying extreme-value copula.

\subsection{Spatial Extremal Dependence via the STANiEXT Framework}

The STANiEXT workflow provides a semiparametric approach for modelling spatial
extremes. It consists of:

\begin{itemize}
\item \textbf{Local GEV fitting:} estimating $(\mu_i,\sigma_i,\xi_i)$ at each grid cell from $\{T_{y,i}\}_y$.
\item \textbf{Dependence estimation:} fitting a max-stable process to obtain pairwise extremal coefficients $\theta_{i,j}$.
\end{itemize}

These coefficients encode the extreme-value copula induced by the fitted spatial model and
form the basis of our clustering approach.

\subsection{Spatial Clustering Based on Extremal Dependence}

We convert extremal coefficients into a dissimilarity matrix:
\[
D_{i,j} = \theta_{i,j} - 1.
\]
Hierarchical clustering of $D$ yields $K$ spatial clusters
\[
C_1,\dots,C_K,
\quad\text{with labels }\ c(i) \in \{1,\dots,K\}.
\]
Clusters represent regions where extreme heat events tend to co-occur and form coherent
units for dependence-respecting loss aggregation.

% ----------------------------------------------------------------------
\section{Methods}

This section describes the steps used to translate spatial extreme temperatures and
population exposure into risk metrics at the grid-cell and cluster levels. The approach
combines local threshold exceedances, loss modelling, and tail-risk via
Value-at-Risk and Expected Shortfall. 

\subsection{Local Thresholds and Exceedances}

Extreme heat is defined relative to the local climate. Following standard practice in the
peaks-over-threshold and climate-extremes literature \citep{coles2001introduction,
huser2020spatial}, we use a high empirical quantile to identify locally extreme events.
For each grid cell $i$, we compute
\[
q_i = \operatorname{quantile}_{0.95}(T_{y,i}).
\]

Given $q_i$, we define the binary exceedance indicator
\[
I_{y,i} = \mathbf{1}(T_{y,i} \ge q_i),
\]
and the continuous exceedance magnitude
\[
S_{y,i} = \max \{ T_{y,i} - q_i,\ 0\}.
\]
Threshold exceedances are widely used in climate-impact modelling because they adapt to
local climatology and avoid arbitrary absolute thresholds
\citep{perkins2013measuring, zhang2011indices}.

\subsection{Population Exposure and Loss Construction}

Population exposure $E_i$ is obtained from high-resolution population datasets and
regridded to match the climate grid using conservative spatial remapping
\citep{schneider2017gpw, bondarenko2020ghsl}. Losses combine exceedance severity with
population exposure:
\[
L_{y,i} = s(T_{y,i}, q_i)\, E_i.
\]

Severity functions are commonly used in climate-damage and environmental-risk modelling
\citep{hallegatte2009future}. We consider:

\begin{itemize}
    \item \textbf{Binary severity}  
    $s(T_{y,i},q_i) = I_{y,i}$, corresponding to event occurrence.

    \item \textbf{Continuous severity}  
    $s(T_{y,i},q_i) = S_{y,i}$, capturing heat intensity beyond the local threshold.
\end{itemize}

\subsection{Risk Measures}

Let $\{L_{y,i}\}_{y=1}^n$ denote the empirical loss sample at grid cell $i$. We use two
standard tail-risk measures extensively applied in finance, insurance, and climate-risk
research \citep{acerbi2002coherence, mcneil2015quantitative}.

The empirical Value-at-Risk (VaR) at confidence level $\alpha$ is
\[
\VaR_{i,\alpha} = \inf \{ \ell : \widehat{F}_{i}(\ell) \ge \alpha \}.
\]

The Expected Shortfall (ES), also known as conditional tail expectation, is
\[
\ES_{i,\alpha}
=
\frac{1}{n(1-\alpha)} \sum_{y : L_{y,i} \ge \VaR_{i,\alpha}} L_{y,i}.
\]

VaR reflects a quantile-level “worst-case threshold”, while ES captures the average
severity of losses conditional on exceeding this threshold. ES is coherent and more
informative about tail behaviour \citep{acerbi2002coherence}.

\subsection{Cluster-Level Aggregation}

Spatial dependence must be accounted for when aggregating climate risks
\citep{koch2020spatialrisk, cooley2019extreme}. Let $C_1,\dots,C_K$ denote clusters
derived from STANiEXT extremal coefficients. These clusters represent coherent regions
where extreme heat events tend to co-occur.

Cluster-level annual losses are computed as
\[
L^{\mathrm{cl}}_{y,k} = \sum_{i \in C_k} L_{y,i}.
\]

Cluster-level risk metrics follow from the empirical distribution of
$\{L^{\mathrm{cl}}_{y,k}\}_y$:
\[
\VaR_{k,\alpha} = \inf \{\ell : \widehat{F}_k(\ell) \ge \alpha\},
\qquad
\ES_{k,\alpha} = \mathbb{E}[L^{\mathrm{cl}}_{y,k} \mid L^{\mathrm{cl}}_{y,k} \ge \VaR_{k,\alpha} ].
\]

Dependence-consistent aggregation ensures that losses from cells that jointly experience
extremes are combined, while independent regions contribute separately
\citep{koch2020spatialrisk}.

\subsection{Interpretation and Use in the Pipeline}

The threshold–severity–exposure framework, combined with spatial dependence-aware
aggregation, provides a robust basis for quantifying extreme-heat risk at both local and
regional scales. The resulting sets of \(\VaR\) and \(\ES\) metrics inform spatial maps,
cluster summaries, and regional comparisons presented in subsequent sections.

































% ----------------------------------------------------------------------
\subsection{Algorithm: Heat-Risk Pipeline}

\begin{tcolorbox}[
    enhanced, breakable,
    colback=white,
    colframe=black,
    title=\textbf{Algorithm 1: Heat-Risk Pipeline},
    fonttitle=\bfseries,
    left=2mm, right=2mm, top=2mm, bottom=2mm, boxrule=0.7pt]

\begin{algorithmic}[1]

\State \textbf{Input:} daily Tmax fields, exposure raster $E_i$, extremal coefficients $\theta_{i,j}$, cluster labels $c(i)$, level $\alpha=0.95$.

\State \textbf{Step 1: Annual maxima}
\[
T_{y,i} = \max_{\text{day}} \{\mathrm{Tmax}_{\text{day},i}\}.
\]

\State \textbf{Step 2: Marginal modelling}
Estimate GEV$(\mu_i,\sigma_i,\xi_i)$ at each grid cell.

\State \textbf{Step 3: Spatial dependence (STANiEXT)}
Retrieve $\theta_{i,j}$ and construct clusters $C_k$ via $D_{i,j} = \theta_{i,j}-1$.

\State \textbf{Step 4: Local thresholds}
\[
q_i = \mathrm{quantile}_{0.95}(T_{y,i}).
\]

\State \textbf{Step 5: Exposure alignment}
Regrid population raster to obtain $E_i$.

\State \textbf{Step 6: Loss computation}
\[
L_{y,i} = s(T_{y,i},q_i)\,E_i.
\]

\State \textbf{Step 7: Cell-wise risk}
Compute $\VaR_{i,\alpha}$ and $\ES_{i,\alpha}$.

\State \textbf{Step 8: Cluster risk}
\[
L^{\mathrm{cl}}_{y,k} = \sum_{i \in C_k} L_{y,i}.
\]
Compute $\VaR_{k,\alpha}$ and $\ES_{k,\alpha}$.

\end{algorithmic}
\end{tcolorbox}

% ----------------------------------------------------------------------
\section{Discussion}

This framework links the local behaviour of extreme heat with its spatial co-movement and
population exposure. Marginal GEV fits quantify location-specific tail behaviour, while
STANiEXT-derived extremal coefficients characterise how extremes propagate across space.
By clustering locations with similar extremal dependence, we construct coherent regions for
dependence-aware aggregation. Within these clusters, Value-at-Risk and Expected Shortfall
provide interpretable metrics of extreme population exposure under rare but plausible heat
events. The resulting maps and cluster summaries highlight both local hotspots and regional
patterns shaped by spatial dependence, offering a comprehensive basis for climate-risk
assessment.











\section{Results}\label{sec:results}

% Ensure nothing from previous sections jumps over this heading
\FloatBarrier

\subsection{World Maps of Tail Risk}\label{sec:maps}

% Figures are forced to appear HERE with [H]
\begin{figure}[H]
  \centering
  \includegraphics[width=\linewidth]{VaR_map.png}
  \caption{Cell-wise $\VaR_{0.95}$ (people). Longitudes wrapped to $[-180,180)$. Spatial dependence not accounted for in this variant.}
  \label{fig:var-map-piecewise}
\end{figure}

\begin{figure}[H]
  \centering
  \includegraphics[width=\linewidth]{ES_map.png}
  \caption{Cell-wise $\ES_{0.95}$ (people). Longitudes wrapped to $[-180,180)$. Spatial dependence not accounted for in this variant.}
  \label{fig:es-map-piecewise}
\end{figure}

\begin{figure}[H]
  \centering
  \includegraphics[width=\linewidth]{VaR_map_optionB.png}
  \caption{Cell-wise $\VaR_{0.95}$ (people) with logarithmic color scale to enhance contrast.}
  \label{fig:var-map-optionB}
\end{figure}

\begin{figure}[H]
  \centering
  \includegraphics[width=\linewidth]{ES_map_optionB.png}
  \caption{Cell-wise $\ES_{0.95}$ (people).}
  \label{fig:es-map-optionB}
\end{figure}

\FloatBarrier


\begin{thebibliography}{99}

\bibitem[Nelsen(2006)]{nelsen2006introduction}
Nelsen, R.~B. (2006).
\newblock \emph{An Introduction to Copulas}.
\newblock Springer Science \& Business Media.

\bibitem[Salvadori et~al.(2007)]{salvadori2007extremes}
Salvadori, G., De~Michele, C., Kottegoda, N.~T., \& Rosso, R. (2007).
\newblock \emph{Extremes in Nature: An Approach Using Copulas}.
\newblock Springer.

\bibitem[Koch and Nolde(2020)]{koch2020spatialrisk}
Koch, E., \& Nolde, N. (2020).
\newblock Spatial risk measures: Local specification and boundary risk.
\newblock \emph{Insurance: Mathematics and Economics}, 92, 87--100.

\bibitem[Davison et~al.(2012)]{davison2012statistical}
Davison, A.~C., Padoan, S.~A., \& Ribatet, M. (2012).
\newblock Statistical modeling of spatial extremes.
\newblock \emph{Statistical Science}, 27(2), 161--186.

\bibitem[Huser and Davison(2014)]{huser2014space}
Huser, R., \& Davison, A.~C. (2014).
\newblock Space--time modelling of extreme events.
\newblock \emph{Journal of the Royal Statistical Society: Series B}, 76(2), 439--461.

\bibitem[Cooley et~al.(2019)]{cooley2019extreme}
Cooley, D., Thibaud, E., \& Davison, A.~C. (2019).
\newblock Extreme value analysis for the environmental sciences.
\newblock \emph{Annual Review of Statistics and Its Application}, 7, 1--26.

\bibitem[Reich and Shaby(2012)]{reich2012hierarchical}
Reich, B.~J., \& Shaby, B.~A. (2012).
\newblock A hierarchical max-stable spatial model for extreme precipitation.
\newblock \emph{Annals of Applied Statistics}, 6(4), 1430--1451.

\bibitem[Coles(2001)]{coles2001introduction}
Coles, S. (2001).
\newblock \emph{An Introduction to Statistical Modeling of Extreme Values}.
\newblock Springer, London.

\bibitem[Davison et~al.(2012)]{davison2012statistical}
Davison, A.~C., Padoan, S.~A., \& Ribatet, M. (2012).
\newblock Statistical modeling of spatial extremes.
\newblock \emph{Statistical Science}, 27(2), 161--186.

\bibitem[Huser and Davison(2020)]{huser2020advances}
Huser, R., \& Davison, A.~C. (2020).
\newblock Advances in statistical modelling of spatial extremes.
\newblock \emph{Wiley Interdisciplinary Reviews: Computational Statistics}, 12(5), e1517.

\bibitem[Huser and Davison(2014)]{huser2014space}
Huser, R., \& Davison, A.~C. (2014).
\newblock Space--time modelling of extreme events.
\newblock \emph{Journal of the Royal Statistical Society: Series B}, 76(2), 439--461.

\bibitem[Cooley et~al.(2019)]{cooley2019extreme}
Cooley, D., Thibaud, E., \& Davison, A.~C. (2019).
\newblock Extreme value analysis for the environmental sciences.
\newblock \emph{Annual Review of Statistics and Its Application}, 7, 1--26.

\bibitem[Reich and Shaby(2012)]{reich2012hierarchical}
Reich, B.~J., \& Shaby, B.~A. (2012).
\newblock A hierarchical max-stable spatial model for extreme precipitation.
\newblock \emph{Annals of Applied Statistics}, 6(4), 1430--1451.

\bibitem[Perkins and Alexander(2013)]{perkins2013measuring}
Perkins, S.~E., \& Alexander, L.~V. (2013).
\newblock Measuring a global set of climate extremes.
\newblock \emph{Journal of Climate}, 26(13), 4500--4519.

\bibitem[Zhang et~al.(2011)]{zhang2011indices}
Zhang, X., Alexander, L., Hegerl, G., Jones, P. et al. (2011).
\newblock Indices for monitoring changes in extremes based on daily temperature and precipitation data.
\newblock \emph{Wiley Interdisciplinary Reviews: Climate Change}, 2(6), 851--870.

\bibitem[Schneider et~al.(2017)]{schneider2017gpw}
Schneider, A., Friedl, M.~A., \& Potere, D. (2017).
\newblock Mapping global urban population using GRUMP and GPW.
\newblock \emph{Remote Sensing of Environment}, 117, 254--267.

\bibitem[Bondarenko et~al.(2020)]{bondarenko2020ghsl}
Bondarenko, M., Kerr, D., \& Sorichetta, A. (2020).
\newblock Estimating population distribution using GHSL.
\newblock \emph{Scientific Data}, 7, 1--12.

\bibitem[Hallegatte et~al.(2009)]{hallegatte2009future}
Hallegatte, S., Hourcade, J.~C., \& Dumas, P. (2009).
\newblock Future climate change impacts on damages from weather-related extremes.
\newblock \emph{Environmental Science \& Policy}, 12(7), 741--758.

\bibitem[Acerbi and Tasche(2002)]{acerbi2002coherence}
Acerbi, C., \& Tasche, D. (2002).
\newblock On the coherence of Expected Shortfall.
\newblock \emph{Journal of Banking \& Finance}, 26(7), 1487--1503.

\bibitem[McNeil et~al.(2015)]{mcneil2015quantitative}
McNeil, A.~J., Frey, R., \& Embrechts, P. (2015).
\newblock \emph{Quantitative Risk Management}.
\newblock Princeton University Press.

\bibitem[Huser and Wadsworth(2020)]{huser2020spatial}
Huser, R., \& Wadsworth, J.~L. (2020).
\newblock Modeling spatial extremes.
\newblock \emph{Annual Review of Statistics and Its Application}, 7, 89--118.

\bibitem[Cooley et~al.(2019)]{cooley2019extreme}
Cooley, D., Thibaud, E., \& Davison, A.~C. (2019).
\newblock Extreme value analysis for the environmental sciences.
\newblock \emph{Annual Review of Statistics and Its Application}, 7, 1--26.

\bibitem[Koch and Nolde(2020)]{koch2020spatialrisk}
Koch, E., \& Nolde, N. (2020).
\newblock Spatial risk measures: Local specification and boundary risk.
\newblock \emph{Insurance: Mathematics and Economics}, 92, 87--100.
























\end{thebibliography}





\appendix
\section{Data Preprocessing Details}
% [Additional diagnostics, masks, lon-wrapping, and plotting settings.]

\section{Parameter Choices}
% [Choices for $\alpha$, severity scale $\Delta$, cluster aggregation (sum vs mean), etc.]







\end{document}


notes for the meeting 18.11.2025